\documentclass[11pt]{article}

\usepackage{amsmath,amssymb,amsthm}
\usepackage{hyperref}
\usepackage[margin=1in]{geometry}

\newcommand{\area}{\operatorname{area}}
\newcommand{\dinv}{\operatorname{dinv}}
\newcommand{\defc}{\operatorname{defc}}
\newcommand{\NRCM}{\operatorname{NRCM}}

\newtheorem{theorem}{Theorem}
\newtheorem{conjecture}{Conjecture}
\newtheorem{remark}{Remark}

\title{Dyck Skeleton String Decompositions}
\author{}
\date{}

\begin{document}

\maketitle

These notes deal with proven and conjectural symmetric string-decomposition
formulas for the \(q,t\)-Catalan polynomials, much of which is built off the
work in \cite{Hawkes2026} and \cite{Hawkes2022}.  In this note, a string means
a sum of monomials in \(q\) and \(t\) in which the powers of \(q\) and \(t\)
respectively increase and decrease step-by-step, and whose total contribution
is symmetric in \(q\) and \(t\).  The status of such decompositions varies
depending on the level of generality: we survey the classical case, the
\(r=\tau s+1\) rational case, and the full rational case.  Within these cases,
we identify important deficit cutoffs and spell out what results are known or
conjectured at those levels.

Two current workspace proof packages substantially refine the earlier status
survey.  The directory
\path{mod1_decomp_formula/} proves the lower-cutoff decomposition and
partition formula for \(r=\tau s+1\), while
\path{naive_rational_cyclic_map.tex} formulates and argues for deficit
preservation of the naive rational cyclic map in the full rational setting.
We distinguish below between the theorem statements in those manuscripts,
their external inputs, and any proof issues still recorded by the accompanying
audit files.

More explicitly, if all terms in a deficit layer have total degree \(N\), then
a positive symmetric string rooted at area \(A\) has the form
\[
  \sum_{j=A}^{N-A} q^j t^{N-j},
\]
provided \(A\le N-A\).  In path language, the terms in such a string have fixed
deficit, area increasing by \(1\), and dinv decreasing by \(1\).  The weak lower
half of a deficit layer means the part where
\[
  \area \le \dinv.
\]
One caution will be useful throughout: a lower-half decomposition proves the
corresponding root formula on the lower half.  It does not by itself prove
\(q,t\)-symmetry or determine the reflected upper half.  For that one needs an
independent symmetry theorem, as in the classical case, or a direct full
coefficient comparison.

\section{The classical Dyck case}

First we summarize the results for the classical case.  An ordinary Dyck
sequence of length \(n\) is a sequence
\[
  x=(x_0,\ldots,x_{n-1})
\]
such that
\[
  x_0=0,\qquad x_i\ge 0,\qquad x_{i+1}\le x_i+1.
\]
We use the usual statistics
\[
  \area(x)=\sum_i x_i
\]
and
\[
  \dinv(x)=
  \#\{(i,j):0\le i<j<n,\ x_i=x_j
  \text{ or } x_i=x_j+1\}.
\]
The top total degree is
\[
  M=\binom n2,
\]
and the deficit is
\[
  \defc(x)=M-\area(x)-\dinv(x).
\]
Thus paths of fixed deficit \(d\) all contribute monomials of the same total
degree \(M-d\).

The classical result uses special Dyck skeletons as the roots of the strings.
A full Dyck skeleton is a Dyck sequence with no extractable entry, using the
leftmost extraction convention from \cite{Hawkes2026}.  For \(n\ge4\), one
exceptional full skeleton is removed:
\[
  \epsilon_n=(0,0,1,\underbrace{0,\ldots,0}_{n-4\text{ entries}},1).
\]
A special Dyck skeleton is a full Dyck skeleton other than \(\epsilon_n\).  For
\(n<4\), every full skeleton is special.

The relevant upper cutoff in the classical case is
\[
  \defc \le 2n-8.
\]
In this range the skeleton-string formula is a theorem.

\begin{theorem}[Hawkes, low-deficit skeleton strings]
For \(n\ge4\) and \(d\le 2n-8\), the weak lower half of the deficit-\(d\)
layer of ordinary Dyck sequences of length \(n\) is partitioned into strings
beginning at special Dyck skeletons.  Along each string, the up map raises area
by \(1\), lowers dinv by \(1\), and preserves deficit.
\end{theorem}

Combining this lower-half decomposition with the known \(q,t\)-symmetry of
\(C_n(q,t)\) gives the special-skeleton formula in the whole low-deficit range:
\[
\left.C_n(q,t)\right|_{\binom n2-2n+8\le \deg_{q,t}\le \binom n2}
=
\sum_{\substack{S\text{ special Dyck skeleton of length }n\\
                \defc(S)\le 2n-8}}
\frac{
q^{\dinv(S)+1}t^{\area(S)}
-
q^{\area(S)}t^{\dinv(S)+1}
}{q-t}.
\]
Equivalently, a skeleton \(S\) of deficit \(d\) contributes the symmetric
interval
\[
  \sum_{j=\area(S)}^{M-d-\area(S)}
  q^j t^{M-d-j}.
\]

In this classical setting, the upper cutoff \(2n-8\) is already a proved
nonnegative range for string decomposition.  The smaller cutoff
\[
  \defc \le n-3
\]
is still important, but for a different reason: in that lower range there is a
simpler partition-indexed description of the roots.  However, the classical
roots in the full range \(\defc\le 2n-8\) are special Dyck skeletons, not simply
all Dyck sequences beginning \((0,0,\ldots)\).

We will not review the local East/West proof of the classical up and down maps
here.  That machinery belongs to the proof in \cite{Hawkes2026}.

\section{The family \(r=\tau s+1\)}

Next, we record the corresponding picture for the rational family
\[
  r=\tau s+1.
\]
This is the family studied in \cite{Hawkes2022} under the notation
\(s=\ell+1\) and \(r=(\ell+1)m+1\), with \(m=\tau\).  The paper
\cite{Hawkes2022} is the source of the broader conjectural rational
string-root formula.  The special-skeleton formula discussed here should be
viewed as a more refined formulation of the same general phenomenon rather than
as a theorem already proved in that paper.

We use normalized \(\tau\)-Dyck sequences of length \(s\):
\[
  x=(x_0,\ldots,x_{s-1}),\qquad
  x_0=0,\qquad x_i\ge0,\qquad x_{i+1}\le x_i+\tau.
\]
The area statistic is again
\[
  \area(x)=\sum_i x_i.
\]
The rational dinv statistic is
\[
\dinv_\tau(x)=\sum_{i<j}
\begin{cases}
\max(0,x_i+\tau-x_j), & x_i\le x_j,\\
\max(0,x_j+1+\tau-x_i), & x_i>x_j.
\end{cases}
\]
The top degree is
\[
  M_{s,\tau}=\tau\binom{s}{2},
\]
and the deficit is
\[
  \defc_\tau(x)=M_{s,\tau}-\area(x)-\dinv_\tau(x).
\]

The notion of a skeleton is defined by an extraction rule analogous to the
classical one, but with the predecessor interval enlarged by \(\tau\).  An
entry \(x_j=e>0\) is extractable if exactly one earlier entry lies in
\[
  [\max(0,e-\tau),e),
\]
and if deleting \(x_j\) still leaves a normalized \(\tau\)-Dyck sequence.  Thus
if \(0<j<s-1\), one also requires the splice condition
\[
  x_{j+1}\le x_{j-1}+\tau.
\]
The leftmost extractable entry is used.  A full rational skeleton is a
normalized \(\tau\)-Dyck sequence with no extractable entry.  For \(s\ge4\),
the exceptional full skeleton is
\[
  \epsilon_{s,\tau}
  =(0,0,1,\underbrace{0,\ldots,0}_{s-4\text{ entries}},\tau),
\]
and a special rational skeleton is a full rational skeleton other than
\(\epsilon_{s,\tau}\).  For \(s<4\), every full rational skeleton is special.

There are two cutoffs to keep separate.  The upper cutoff is
\[
  B_{\mathrm{upper}}=(s-2)(\tau+1)-4.
\]
In the range
\[
  \defc_\tau\le B_{\mathrm{upper}},
\]
there appears to be a special rational skeleton formula:
\[
\sum_{\substack{
S\text{ special rational skeleton}\\
\defc_\tau(S)\le B_{\mathrm{upper}}
}}
\frac{
q^{\dinv_\tau(S)+1}t^{\area(S)}
-
q^{\area(S)}t^{\dinv_\tau(S)+1}
}{q-t}.
\]
Equivalently, a skeleton \(S\) of deficit \(d\) should contribute the interval
\[
  \sum_{j=\area(S)}^{M_{s,\tau}-d-\area(S)}
  q^j t^{M_{s,\tau}-d-j}.
\]
This is the direct rational analogue of the classical special-skeleton
formula.  At present we regard it as conjectural, supported by coefficient
checks.  In the checked cases, the refined special skeletons in this range give
nonnegative interval contributions.  We do not currently claim that the
upper-cutoff formula comes from a complete up-map decomposition.

The lower cutoff is
\[
  B_{\mathrm{lower}}=(s-2)\tau-1.
\]
In this smaller range, the situation is explicitly decompositional.  The
manuscripts in
\path{mod1_decomp_formula/} now prove the required lower-half decomposition.
The roots of the strings are exactly the full skeletons in the defect range,
and these are exactly the normalized \(\tau\)-Dyck sequences beginning
\[
  (0,0,\ldots).
\]
Equivalently, these are the paths whose label-\(1\) position has value \(0\).
The final formula manuscript proves the stronger root estimate
\[
  \area(S)\le\defc_\tau(S),
\]
so all roots in this range lie strictly below the middle of their defect
layers.

\subsection{The proved lower-cutoff decomposition}

The proof package splits the construction by length:
\begin{itemize}
\item \path{sleq4_decomp.tex} proves the lower-half
  decomposition for \(s=3,4\).  East3 and West3 cannot fail for \(s=3\).
  At \(s=4\), all failed level-three stages belong to the single inverse
  boundary pair
  \[
    (0,p+1,a,c)\longleftrightarrow(0,a+1,c+1,p),
  \]
  with \(0\le p<a\le\tau-1\) and \(p+\tau<c\le a+\tau\).
\item \path{s5_decomp.tex} proves the \(s=5\)
  lower-half decomposition.  It handles the genuinely new failed-West5
  boundary by mutually inverse saving and repacking operations.
\item \path{sgeq6_decomp.tex} proves the uniform
  decomposition for \(s\ge6\), conditional in its statement on the
  special-skeleton position and East5/West5 nonfailure inputs.
\item \path{east5_nonfailure.tex} proves that a failed
  East5 stage in the relevant defect range lies strictly above the midpoint,
  and \path{west5_nonfailure.tex} proves West5 nonfailure
  for the weak-lower-half DOWN procedure when \(s\ge6\).
\item \path{mod1_formula.tex} proves the root-position
  input, assembles the three length ranges, and identifies the roots with
  bounded partitions using the bijection from \cite{Hawkes2022}.
\end{itemize}

The resulting formula can be stated as follows.

\begin{theorem}[Modulo-one lower-cutoff partition formula]
Fix \(\tau\ge2\) and \(s\ge1\), and put
\[
 M_{s,\tau}=\tau\binom{s}{2},
 \qquad B=(s-2)\tau-1.
\]
Assume the known \(q,t\)-symmetry of
\(\operatorname{Cat}_{\tau s+1,s}(q,t)\).  For \(s\ge3\),
\begin{align*}
 &\sum_{\substack{x\text{ normalized}\\\defc_\tau(x)\le B}}
 q^{\dinv_\tau(x)}t^{\area(x)}\\
 &\quad=
 \sum_{\substack{\lambda\text{ a partition}\\
                 \lambda_1\le s-2,\ |\lambda|\le B}}
 \frac{
 q^{M_{s,\tau}-|\lambda|-\ell(\lambda)+1}t^{\ell(\lambda)}
 -
 q^{\ell(\lambda)}
 t^{M_{s,\tau}-|\lambda|-\ell(\lambda)+1}
 }{q-t}.
\end{align*}
For \(s=1,2\), the range \(\defc_\tau\le B\) is empty, so the assertion is
vacuous.  Equivalently, the right side may be indexed by the full skeletons
of defect at most \(B\), or by the normalized words beginning with two zeros
in that range.
\end{theorem}

Thus the two cutoffs now have sharply different statuses.  The refined
special-skeleton formula at \(B_{\mathrm{upper}}\) remains conjectural and
is supported by coefficient checks.  The simpler double-zero and
partition-indexed formula at \(B_{\mathrm{lower}}\) is proved by the
manuscripts in \path{mod1_decomp_formula/}, using the known symmetry and the
Hawkes bounded-partition bijection as external inputs.

\section{The full rational case}

Finally, we discuss the full rational case.  Let \(r\) and \(s\) be coprime
positive integers.  In this setting it is useful to work with position
coordinates.  For \(0\le i<s\), define
\[
  H_i=\left\lfloor \frac{ri}{s}\right\rfloor,
  \qquad
  L_i\equiv ri \pmod s,\qquad 0\le L_i<s.
\]
The numbers \(L_i\) are the labels of the columns.  Since \(r\) and \(s\) are
coprime, every label occurs exactly once.  Thus the phrase ``the label-\(1\)
position'' means the unique index \(i\) such that \(L_i=1\).

A rational Dyck path in these coordinates is a sequence
\[
  Q=(Q_0,\ldots,Q_{s-1})
\]
such that
\[
  Q_0=0,\qquad 0\le Q_i\le H_i,
\]
and whose path heights
\[
  P_i=H_i-Q_i
\]
are weakly increasing:
\[
  P_0\le P_1\le\cdots\le P_{s-1}.
\]
We take
\[
  \area(Q)=\sum_i Q_i.
\]
Let
\[
  M=\sum_i H_i.
\]
For the NRCM checks it is convenient to use the following equivalent deficit
statistic.  For \(1\le i<j<s\), put \(u=|Q_i-Q_j|\).  If \(Q_i\ne Q_j\) and
the comparisons \(Q_i>Q_j\) and \(L_i>L_j\) have opposite truth values, replace
\(u\) by \(u-1\), and then clamp \(u\) to be nonnegative.  Let
\(\Delta_i=H_{i+1}-H_i\), and define
\[
  v_{ij}=\begin{cases}
  \Delta_i-(Q_{i+1}-Q_i), & Q_i>Q_j,\\
  \Delta_{i-1}-(Q_i-Q_{i-1}), & Q_j>Q_i,\\
  0, & Q_i=Q_j.
  \end{cases}
\]
Then
\[
  \defc(Q)=\sum_{1\le i<j<s}\min(u_{ij},v_{ij}),
  \qquad
  \dinv(Q)=M-\area(Q)-\defc(Q).
\]
Here \(v_{ij}\) is not separately clamped; this is the statistic used in the
code checks.

As before, the weak lower half of a deficit layer is the region
\[
  \area(Q)\le \dinv(Q),
\]
or equivalently
\[
  \area(Q)\le \left\lfloor\frac{M-\defc(Q)}2\right\rfloor.
\]

The general rational analogue of beginning \((0,0,\ldots)\) is having value
\(0\) in the label-\(1\) position.  Indeed, in the special family
\(r=\tau s+1\), one has
\[
  L_i\equiv i\pmod s,
\]
so the label-\(1\) position is \(i=1\), and the condition is exactly that the
path begins \((0,0,\ldots)\).

For the full rational case we do not give a simple numerical formula for the
lower cutoff.  Instead, we define the lower cutoff intrinsically: it is the
largest deficit range in which every path with value \(0\) in the label-\(1\)
position lies in the weak lower half.  This definition agrees with the lower
cutoff \(n-3\) in the classical case and with
\[
  (s-2)\tau-1
\]
in the family \(r=\tau s+1\).

In light of the conjectural formula from \cite{Hawkes2022}, this intrinsic
lower cutoff is the natural range in which one should expect a positive
string-decomposition formula.  The reason is that the paths with value \(0\)
in the label-\(1\) position are precisely the predicted roots of the strings.
When all of these roots lie in the weak lower half, the corresponding string
contributions run from the root up to its symmetric partner without introducing
negative terms.  Outside this range, the same rational formula may still make
sense, but it can involve negative terms and should not be interpreted as a
literal positive string decomposition without additional structure.

\subsection{The naive rational cyclic map}

The naive rational cyclic map, or NRCM, is one piece of such additional
structure.  Given a path \(Q\), and for each \(k\ge1\), let
\[
  I_k=\{k,k+1,\ldots,s-1\}.
\]
List the indices in \(I_k\) in increasing label order.  The candidate move
\(T_k(Q)\) cyclically moves the \(Q\)-values along this label-ordered suffix:
each value moves to the next suffix position in label order, and the last value
wraps around to the first suffix position after being increased by \(1\).
Values outside \(I_k\) are unchanged.

The strict NRCM chooses the first \(k\), starting from \(k=1\), for which this
candidate \(T_k(Q)\) satisfies the capacity inequalities
\[
  0\le T_k(Q)_i\le H_i.
\]
If this first capacity-valid candidate is also a valid rational Dyck path, then
\[
  \NRCM(Q)=T_k(Q).
\]
If the first capacity-valid candidate is not path-valid, then NRCM is
undefined.

By construction, every defined NRCM move raises area by \(1\).  The main
result stated in \path{naive_rational_cyclic_map.tex} is
\[
  \defc(\NRCM(Q))=\defc(Q)
\]
whenever \(Q\) and \(\NRCM(Q)\) are valid and the map is defined.  Consequently
a defined NRCM move would change the three string statistics by
\[
 (\area,\dinv,\defc)\longmapsto
 (\area+1,\dinv-1,\defc).
\]

The logical status of this statement should be recorded precisely.  The
manuscript gives a detailed threshold-count and endpoint-cancellation proof,
but the workspace audit \path{NRCM_report.md}, accepted in
\path{NRCM_reply.md}, identifies two repairs still needed in that proof: the
transport formulas do not cover a one-column selected suffix, and the central
endpoint-pairing argument has not yet been derived at the required token-level
precision.  Thus deficit preservation is the principal NRCM theorem supported
by the manuscript and extensive finite checks, but within the current
workspace it remains pending those proof repairs rather than being used as an
unconditional input to the modulo-one theorem above.

The paths with value \(0\) in the label-\(1\) position naturally appear as
NRCM bases: such a path cannot be the image of an NRCM move.  Informally, the
inverse cyclic motion would require undoing a wraparound increase at the first
available label, and a zero in the label-\(1\) position gives no room for that
inverse move.  This agrees with the roots predicted by the rational
\(q,t\)-Catalan conjecture.

Even after deficit preservation is fully repaired, NRCM should not be
confused with a complete proof of the intrinsic lower-cutoff decomposition.
In general, NRCM does not give a full lower-half decomposition throughout that
entire cutoff.  The code checks the label-\(1\)-zero positivity prediction
separately from NRCM coverage.  The computations show genuine NRCM
lower-half decompositions only in certain smaller, case-dependent deficit
ranges.  In the family \(r=\tau s+1\), whenever both maps are defined, NRCM
agrees with the proved modulo-one UP map described in
\path{mod1_decomp_formula/}.

\section{Computational checks}

The accompanying file \texttt{code.py} is a consolidated checker for the claims
discussed here.  It separates three logically different tests.

First, it performs finite sanity checks of the classical special-skeleton
coefficient formula.  These checks are not needed for the theorem, but they
verify that the definitions in the code agree with the stated formula in small
cases.

Second, for \(r=\tau s+1\), it separates the two rational cutoffs.  At the
upper cutoff
\[
  (s-2)(\tau+1)-4,
\]
it compares direct coefficient dictionaries with the special-rational-skeleton
formula.  At the lower cutoff
\[
  (s-2)\tau-1,
\]
it checks that the special skeleton roots agree with the paths beginning
\((0,0,\ldots)\), that these roots lie in the weak lower half, that the
corresponding positive interval formula matches direct enumeration, and that
the East3/East5 up and down maps give lower-half coverage in the checked
nondegenerate cases.  The code deliberately does not use upper-cutoff East7
obstructions as evidence against the lower-cutoff decomposition claim.
These lower-cutoff computations are now regression checks for the proof in
\path{mod1_decomp_formula/}, rather than the primary evidence for that claim.

The original item-level \(r=\tau s+1\) runner is configured on the grid
\[
\begin{array}{c|c}
\tau & s\text{-range}\\
\hline
2 & 1\le s\le14\\
3 & 1\le s\le12\\
4 & 1\le s\le10\\
5 & 1\le s\le9.
\end{array}
\]
The consolidated \texttt{code.py} includes an \texttt{official-r1mod} preset
following this grid, but with the interpretation used in this note: formula
checks belong to the upper cutoff, while map/decomposition checks belong to the
lower cutoff.  The separate manuscripts, rather than this finite grid, prove
the lower-cutoff result for all \(\tau\ge2\) and all lengths.

Third, in the full rational case, the code computes the intrinsic lower cutoff,
checks the positive label-\(1\)-zero interval formula through that cutoff, and
separately checks NRCM properties and NRCM lower-half coverage.  The latter is
expected to work only in smaller case-by-case ranges, not throughout the full
intrinsic lower cutoff.  These computations also support the
deficit-preservation statement in \path{naive_rational_cyclic_map.tex}, but do
not replace the two proof repairs recorded in \path{NRCM_report.md}.

\section{Summary of cutoffs}

The three cases can be summarized as follows.

\begin{center}
\small
\renewcommand{\arraystretch}{1.25}
\begin{tabular}{p{0.20\textwidth}|p{0.18\textwidth}|p{0.22\textwidth}|p{0.28\textwidth}}
\textbf{case} & \textbf{upper cutoff} & \textbf{lower cutoff} & \textbf{status} \\
\hline
classical Dyck
&
\(2n-8\)
&
\(n-3\)
&
proved special-skeleton decomposition at upper cutoff
\\
\(r=\tau s+1\)
&
\((s-2)(\tau+1)-4\)
&
\((s-2)\tau-1\)
&
upper formula conjectural; lower decomposition and partition formula proved
\\
general rational
&
no NRCM upper cutoff
&
intrinsic label-\(1\)-zero lower cutoff
&
positive formula conjectural; NRCM preservation proof pending recorded repairs
\end{tabular}
\end{center}

The important distinction is that the upper cutoff and lower cutoff play
different roles.  In the classical case, the upper cutoff already gives a
proved nonnegative special-skeleton decomposition.  In the \(r=\tau s+1\)
case, the upper cutoff gives a conjectural nonnegative refined-skeleton
formula, while the lower cutoff has a proved lower-half decomposition with the
simpler roots \((0,0,\ldots)\) and a bounded-partition formula.  In the general
rational case, the lower cutoff is defined by the condition that all
label-\(1\)-zero paths lie in the weak lower half.  NRCM proposes a structural
move with the correct statistic change and gives complete strings in smaller
examples, but it does not cover the whole lower-cutoff decomposition in
general, and its current deficit-preservation proof still requires the repairs
recorded above.

\begin{thebibliography}{9}

\bibitem{Hawkes2026}
Graham Hawkes.
\newblock \emph{Dyck Symmetric Functions and Applications to
\(q,t\)-Catalan Polynomials}.
\newblock arXiv:2605.13003, 2026.
\newblock \url{https://arxiv.org/abs/2605.13003}.

\bibitem{Hawkes2022}
Graham Hawkes.
\newblock \emph{A conjectured formula for the rational \(q,t\)-Catalan
polynomial}.
\newblock arXiv:2208.00577, 2022.
\newblock \url{https://arxiv.org/abs/2208.00577}.

\end{thebibliography}

\end{document}
